\chapter{Заключение Тома VI: различающий тест рождения материи}

Том VI был открыт не для классификации ещё одного допустимого дефекта, а
для проверки более сильного утверждения: способен ли один физически
определённый родительский закон сам вывести нулевой вакуум из равновесия и
нелинейно завершить переход устойчивой глобально нейтральной Real-парой.
Рабочая гипотеза допускала такой процесс; нулевая гипотеза требовала
остановить программу, если общий родитель, запуск или устойчивое
насыщение не будут получены.

Различающий тест выполнен. Том существенно расширил строгий
математический слой, но рабочая динамическая гипотеза в текущей
архитектуре не подтверждена. Это не универсальная теорема невозможности
рождения материи. Это точный отрицательный результат для объявленного
пространства носителей, действий, связей и дискретных продолжений.

\section{Ответ на центральный вопрос}

Вопрос Тома VI имел форму
\begin{equation}
 I\longrightarrow T_{\rm sing}
 \longrightarrow(V_{15},V_{-15}),
 \label{eq:v6-final-central-process}
\end{equation}
где все три стадии должны принадлежать одной конфигурационной задаче и
одному родительскому функционалу.

Получен частичный общий операторный язык: индекс-нулевая пара допускает
конечноранговый обменный мост, ориентированные ветви имеют классы
$+15/-15$, а смена ранга и спектральный поток типизированы. Однако
единого физического действия, которое одновременно содержит нулевой
вакуум, проекторный параметр порядка, локализованный конечный объект и
исходящий канал, не построено. Поэтому процесс
\eqref{eq:v6-final-central-process} не выведен как автономная динамика.

\begin{theorem}[Ответ Тома VI]
В проверенной архитектуре не доказано самопроизвольное рождение материи
из нулевого вакуума. Классификация Real-пары, кинематика переходного слоя,
статический проекторный фазовый переход и несколько точных
локализованных решений замкнуты, но единый родитель, эндогенный запуск,
физическая скорость и устойчивый наблюдаемый конечный объект отсутствуют.
Следовательно, рабочая динамическая гипотеза не прошла критерий успеха
Тома VI, а продолжение той же архитектуры остановлено.
\end{theorem}

\section{Сопоставление исходных задач и результата}

\begin{center}
\footnotesize
\setlength{\tabcolsep}{3pt}
\begin{longtable}{>{\raggedright\arraybackslash}p{0.08\textwidth}>{\raggedright\arraybackslash}p{0.25\textwidth}>{\raggedright\arraybackslash}p{0.48\textwidth}>{\raggedright\arraybackslash}p{0.12\textwidth}}
\toprule
№ & Исходная задача & Фактический результат & Статус \\
\midrule
\endhead
1 & Общее пространство $I$, $T_{\rm sing}$ и Real-пары & Обменное
операторное пространство построено; единый физический носитель всех
динамических секторов не атрибутирован & частично \\
2 & Типизация потери обратимости & Индекс, ядро, коядро, clutching и
Callias--Toeplitz-сравнение замкнуты кинематически & пройдено \\
3 & Один родительский функционал & Проекторная энергия, compacton-правило,
граница, аффинная связь и радиация происходят из разных конструкций & не
пройдено \\
4 & Физический гессиан нулевого вакуума & Вычислены многочисленные
секторные гессианы, но общего физического гессиана одного родителя нет &
не пройдено \\
5 & Нелинейное насыщение Real-парой & Найден точный compacton, но он не
имеет вакуумного бассейна захвата и не является физической Real-парой
$+15/-15$ & не пройдено \\
6 & Конечные энергия и радиус & Получены безразмерные профили, натяжение
прямой нити и условные масштабные произведения; абсолютный устойчивый
радиус не выведен & частично \\
7 & Полная устойчивость & Для бозонной прямой нити закрыта поперечная
линейная устойчивость; замкнутое кольцо и compacton как частица не
стабилизированы & частично \\
8 & Квантовая мера и Pfaffian-линия & Real-суммирование и эта-фазы
проверены; затухание и селекция ветви не следуют из них & частично \\
9 & Карта к массам и взаимодействиям & Порталы, абсолютный масштаб,
массы, времена жизни и плотности не выведены; феноменологическая карта
остановлена & не пройдено \\
10 & Слепые наблюдательные тесты & Динамическая ветвь не дала независимого
размерного или безразмерного наблюдаемого следствия, не использованного
как вход & не пройдено \\
\bottomrule
\end{longtable}
\end{center}

Таким образом, Том VI достиг цели как различающий исследовательский
протокол, но не достиг заявленного положительного физического замыкания.

\section{Главные положительные результаты}

\subsection{Обменная Real-пара и кинематика ранга}

Различены три понятия: ориентированный индекс, обратимость конкретного
представителя и стабильный Real-класс. Индекс-нулевая пара допускает
обратимое конечноранговое дополнение, тогда как каждая ориентированная
ветвь сохраняет препятствие ранга пятнадцать. Поэтому дефицит $1/7$ не
может одновременно читаться как неизбежная размерность ядра любого
представителя и как нормированный стабильный класс.

Clutching-символ проекторного spin-cover и тёплицев символ сопоставлены на
уровне комплексного и Real-парного классов. Класс пятнадцать закреплён как
реестр полного коэффициентного пакета одного поколения, но не как число
частиц одного события.

\subsection{Статический проекторный переход}

Для вещественного семейного состояния построен точный контрольный
фазовый ландшафт. Канонический нормированный функционал имеет переход
первого рода с
\begin{align}
 \beta_c&=1.5426695408602848\ldots,\\
 \operatorname{Spec}R_c
 &=(0.912166596\ldots,0.043916702\ldots,
 0.043916702\ldots),
\end{align}
а изотропная ветвь теряет локальную устойчивость при
\begin{equation}
 \beta_{\rm sp}=\frac{21}{2}.
\end{equation}
Поле
\begin{equation}
 Q=R-\frac13I_3
\end{equation}
имеет пять вещественных мод: одну амплитудную, две биаксиальные и две
директорные. Его вакуумная орбита равна $\mathbb{RP}^2$, а топологические
сектора включают линейные, точечные и хопфовы конфигурации.

Энтропийный и энергетический бюджеты перехода вычислены:
\begin{equation}
 \Delta S_R=0.7402345697629901\ldots,
 \qquad
 \Delta E_R=0.479840011199152\ldots.
\end{equation}
Доказана кинематическая вместимость конечного внутреннего носителя, но не
необратимый закон охлаждения.

\subsection{Бозонный дефект и струна}

Сохранившийся объект классифицирован как нейтральный бозонный семейный
квинтет, а не как наблюдаемая фермионная частица. Для связанного
$Q$-тензорного и калибровочного родителя построены вакуум, вихревой
профиль и полный тензорный оператор возмущений. После удаления
переносных мод получена положительная внутренняя щель и закрыта поперечная
линейная устойчивость прямой нити.

Из натяжения выведен член Намбу--Гото эффективного действия. Однако
контролируемая жёсткость не стабилизирует замкнутое кольцо, нормируемый
ток мирового листа не получен, а радиус основного кольца стремится к
нулю. Поэтому прямая устойчивая нить не повышена до частицы конечной
массы.

\subsection{Геометрическая нить и спектральный примитив}

Многопроходное чтение моментов и абстрактный незаветвляющийся цикл
математически допустимы. Но конечная толщина, исключённый объём и
бесконечный барьер пересоединения не следуют из родителя. Буквальная
неразрывная нить требует нового микроскопического постулата или масштаба.

Поэтому геометрическая метафора была дисциплинированно заменена
спектральным примитивом: билатеральным сдвигом, оператором числа,
Hardy-проектором и конечным дефектным проектором ранга пятнадцать. Это
сохраняет историю переходов, не объявляя её материальной верёвкой.

\subsection{Точный нелинейный дискретный результат}

На пространственно-хиральном носителе построено локальное
нормосохраняющее нелинейное обновление. Оно имеет точные двухузловые
compacton-решения
\begin{equation}
 \kappa=2(2m+1)\pi,
 \qquad
 \kappa_{\min}=2\pi,
 \qquad
 F^2(\Psi)=-\Psi.
\end{equation}
Вычислен характерно разрешённый исходящий канал
\begin{equation}
 \rho_{\rm rad}(0)=2\pi,
 \qquad
 \Gamma_{\rm walk}=4\pi^2
\end{equation}
за четырёхтактный цикл на квадрат касательной амплитуды. Этот коэффициент
переживает полуслед полной Real-пары.

Одновременно доказаны границы результата: точные решения образуют
непрерывное нейтральное многообразие, обычные локализованные данные не
захватываются, эта/Pfaffian-фаза не создаёт затухание, абсолютный масштаб
остаётся свободным, а факторизованный подъём радиации точно сохраняет
семейное состояние $R$.

\section{Главные отрицательные результаты}

Строгие запреты Тома VI не позволяют повторять без новой архитектуры:
\begin{enumerate}
 \item отождествление обратимого представителя полной пары с нулевым
 Real-классом;
 \item вывод проекторной фазы из линейного функционала Гиббса без
 нелинейного инварианта;
 \item получение самозапуска из канонического обменного моста, который
 стабилизирует изотропию;
 \item выбор ранга обычной полярной или Wishart-мерой;
 \item объявление внутренней унитарной передачи необратимым охлаждением;
 \item использование конечных часов без обратной реакции и рекурренций;
 \item чтение Callias-, сфалеронного или аномального индекса как
 вероятности рождения класса пятнадцать;
 \item объявление класса пятнадцать числом частиц или элементарным
 неделимым объектом;
 \item превращение буквальной многопроходной нити в физический объект без
 толщины и барьера пересоединения;
 \item выбор одной compacton-ветви граничным условием, эта-фазой или
 вставленным вручную коэффициентом затухания;
 \item чтение радиационной нормы $4\pi^2$ как теплового тока без общего
 действия и следовой нормировки;
 \item добавление $Q$- или $P$-зависимого переплетающего оператора после
 того, как требуемая ось уже выбрана.
\end{enumerate}

Эта карта запретов является самостоятельным результатом: она сокращает
пространство следующих моделей и препятствует круговому возвращению к
уже проверенным объяснениям.

\section{Степень достижения цели}

Финальный динамический контракт имеет вид
\begin{equation}
 \boxed{
 R0:\ \text{частично},\quad
 R1\text{--}R5:\ \text{провал},\quad
 R6:\ \text{пройдено}.}
\end{equation}
Типизированные носители существуют по отдельности, но один
взаимодействующий носитель не построен. Единого действия, запуска, меры,
скорости, устойчивого конечного объекта и слепого динамического
предсказания нет. Воспроизводимость и заранее заданные стоп-критерии
выполнены.

Следовательно, Том VI завершён в двух разных смыслах:
\begin{enumerate}
 \item как математическое и классификационное исследование он дал новые
 строгие результаты;
 \item как положительная физическая теория автономного рождения материи
 он не замкнут и заморожен.
\end{enumerate}

\section{Условия следующей программы}

Следующая программа не должна начинаться с нового локального исправления
compacton, выбора другой $C_4$-оси или перенормировки уже найденного
коэффициента. До открытия нового тома требуется предъявить один новый
родительский кандидат, проходящий предварительный входной контроль.

\subsection*{P0: общий физический носитель}

Носитель должен одновременно содержать семейный параметр порядка,
хиральный коэффициентный сектор, пространственные моды и Real-структуру.
Прямое тензорное произведение без нескалярного взаимодействия не проходит.

\subsection*{P1: одно действие и одна нормировка}

Вакуум, переходный слой, локализованный объект и исходящий канал должны
следовать из одного действия с одним следом, одной мерой и одной системой
размерностей. Секторные коэффициенты после вычисления результата запрещены.

\subsection*{P2: физический запуск}

До анализа конечного объекта необходимо вычислить физический гессиан
нулевого сектора или конечное квантовое действие нуклеации. Запуск должен
существовать на ненулевом множестве вакуумно допустимых данных.

\subsection*{P3: нелинейный конечный объект}

Развитие нестабильности должно завершаться изолированной устойчивой
Real-парой либо другим заранее типизированным наблюдаемым объектом. Седло,
непрерывное нейтральное многообразие или прямая бесконечная нить не
считаются конечным результатом.

\subsection*{P4: мера и скорость}

Требуются мера траекторий, предэкспоненциальный множитель, физическая
единица времени и проверка сохранения энергии или контролируемого
открытого предела. Топологический индекс и норма исходящего пакета не
заменяют скорость.

\subsection*{P5: масштаб и слепое следствие}

До сравнения с данными должны быть зарегистрированы хотя бы одно
абсолютное либо безразмерное наблюдаемое следствие и сертификат провала.
Масса, радиус или время жизни не могут одновременно служить входом и
предсказанием.

\subsection*{P6: воспроизводимость}

Новый кандидат обязан с первого гейта иметь связанную триаду
``гейт--аудит--результат'', явную нулевую гипотезу и критерий остановки.
Положительные и отрицательные исходы имеют одинаковый статус данных.

\begin{nogocriterion}[Запрет автоматического открытия следующего тома]
Следующий том не открывается одним новым названием. Сначала должен быть
предъявлен родительский кандидат, который по крайней мере одновременно
проходит $P0$ и $P1$ и формулирует вычислимый тест $P2$. До этого
допустимы только синтез, внешний аудит и независимое воспроизведение
закрытого корпуса.
\end{nogocriterion}

\section{Финальная формула статуса}

\begin{equation}
 \boxed{
 \begin{gathered}
 \text{Real-пара и её классы классифицированы,}\\
 \text{статическая проекторная фаза и поле порядка построены,}\\
 \text{новые локализованные и радиационные решения получены,}\\
 \text{автономная динамика рождения материи не выведена.}
 \end{gathered}}
 \label{eq:v6-final-status-box}
\end{equation}

\begin{protocol}[Финальная заморозка Тома VI]
\begin{enumerate}
 \item общий операторный язык: \textbf{построен частично};
 \item Real-классы $+15/-15$: \textbf{сохранены};
 \item статический переход: \textbf{$\beta_c=1.5426695409\ldots$};
 \item спинодаль: \textbf{$\beta_{\rm sp}=21/2$};
 \item поле порядка: \textbf{$Q\in\operatorname{Sym}_0(3,\mathbb R)$};
 \item нейтральный бозонный дефект: \textbf{кандидат, не частица СМ};
 \item прямая бозонная нить: \textbf{линейно устойчива поперечно};
 \item замкнутое кольцо конечного радиуса: \textbf{не получено};
 \item буквальная единая нить ткани: \textbf{не выведена};
 \item точный compacton: \textbf{$\kappa_{\min}=2\pi$};
 \item радиационный коэффициент: \textbf{$4\pi^2$};
 \item радиационное охлаждение $R$: \textbf{закрыто отрицательно};
 \item единый родитель и запуск: \textbf{не получены};
 \item физическая скорость и абсолютный масштаб: \textbf{не получены};
 \item автономное рождение материи: \textbf{не выведено};
 \item динамическая архитектура Тома VI: \textbf{заморожена};
 \item следующий допустимый шаг: \textbf{синтез, внешний аудит или новый
 родитель, проходящий $P0$--$P2$}.
\end{enumerate}
\end{protocol}